% संपूर्ण मराठी ग्रंथातील प्रकरण 59: $\Z$ कडे वाटचाल
% हा संपादनयोग्य स्रोतखंड आहे. संकलनासाठी संपूर्ण स्रोतसंग्रहातील
% अक्षररूपे, पूर्वभाग, चिन्हव्याख्या आणि आकृती-स्रोत आवश्यक आहेत.
% मूळ: Open Logic Project; मजकूर आणि रूपांतर CC BY 4.0.
% यंत्रानुवाद, दुरुस्ती आणि तपासणी: OpenAI Codex —
% GPT-5.6 Sol आणि GPT-6 Sol, दोन्ही Ultra effort.
% दुरुस्ती-पुनरावलोकन: GPT-6.1 Sol, Ultra effort.
\chapter{$\Z$ कडे वाटचाल}\label{sth:z:chap}




\section{मांडणीचे आणखी तपशील}\label{sth:z:story:sec}

\ref{sth:story:approach:sec} मध्ये आपण संच तयार होण्याच्या प्रक्रियेबद्दल शोनफिल्ड यांचे वर्णन उद्धृत केले. आता ही मांडणी थोडी अधिक नेमकी करण्यासाठी आणखी काही तत्त्वे लिहू. ती अशी:
\begin{enumerate}
\item[] \stageshier. प्रत्येक संच कोणत्या ना कोणत्या टप्प्यावर तयार होतो.
\item[] \stagesord. टप्प्यांचा क्रम असतो: काही टप्पे इतरांच्या \emph{आधी} येतात.\footnote{प्रत्यक्षात टप्पे \emph{सुक्रमित} आहेत, असे आपण अप्रकटपणे गृहीत धरू. याचा अर्थ \ref{sth:ordinals:chap} मध्ये स्पष्ट केला आहे. हे महत्त्वाचे गृहीतक आहे. खरे तर Scott (1974) यांच्या अतिशय कल्पक पद्धतीने हे गृहीतक प्रथम \emph{टाळता} येते आणि मग \emph{सिद्ध} करता येते. (आरंभीचा टप्पा असावा असे का मानावे, हेही त्यातून स्पष्ट होईल.) येथे त्यात जाता येणार नाही; अधिक माहितीसाठी Button (2021) पाहा.}
\item[] \stagesacc. कोणताही टप्पा $S$ आणि $S$ च्या \emph{आधी} तयार झालेल्या संचांचा कोणताही संग्रह घेतल्यास, टप्पा~$S$ वर नेमके तेच संच सदस्य असलेला संच तयार होतो. टप्पा~$S$ वर दुसरे काही तयार होत नाही.
\end{enumerate}
ही अनौपचारिक तत्त्वे आहेत; पण त्यांच्या आधारे झर्मेलो यांच्या संच उपपत्तीतील अनेक स्वयंसिद्धकांचे समर्थन करता येईल.

(एक सावधानतेची नोंद. पुढे आपण सर्वमान्य नावे असलेली प्रमाणित स्वयंसिद्धके मांडणार आहोत; पण नुकतीच तिरपी अक्षरे वापरून मांडलेली तत्त्वे साहित्यात कोणत्याही विशिष्ट नावांनी ओळखली जात नाहीत. आपण त्यांना केवळ उपयोगी वाटतील अशी नावे दिली आहेत.)








\section{विभक्तीकरण}\label{sth:z:sep:sec}

अनौपचारिक गुणधर्माधारित संचरचनेच्या जागी पुढील तत्त्व घेऊ:

\begin{axiom}[विभक्तीकरण-रूपबंध] प्रत्येक सूत्र $\phi(x)$ साठी पुढील एक स्वयंसिद्धक आहे: कोणताही $A$ दिल्यास $\Setabs{x \in A}{\phi(x)}$ हा संच अस्तित्वात असतो.
\end{axiom}

हे एकच स्वयंसिद्धक नाही, हे लक्षात घ्या. हा स्वयंसिद्धकांचा \emph{रूपबंध} आहे. विभक्तीकरणाची स्वयंसिद्धके \emph{अनंत} आहेत; प्रत्येक सूत्र $\phi(x)$ साठी एक. हा रूपबंध पुढीलप्रमाणेही लिहिता येतो (आणि सहसा तसाच लिहिला जातो):

\begin{defish}
``$S$'' नसलेल्या कोणत्याही सूत्र $\phi(x)$ साठी पुढील एक स्वयंसिद्धक आहे:
\[
	\forall A \exists S \forall x(x \in S \liff (\phi(x) \land x \in A)).
\]
\end{defish}

\ref{sth:::part} च्या सुरुवातीला सांगितलेल्या संकेतानुसार, विभक्तीकरणाच्या स्वयंसिद्धकांतील सूत्रांना~$\phi$ प्राचले असू शकतात.\footnote{याचा अर्थ काय, यासाठी \ref{sfr:infinite:induction:natinductionschema} नंतरची चर्चा पाहा.}

आपण सांगत आलेली संचांची संचयी-क्रमिक संकल्पना विभक्तीकरणाचे लगेच समर्थन करते. ते पाहण्यासाठी $A$ हा संच मानू. \stageshier{} नुसार $A$ कोणत्या तरी टप्प्या~$S$ वर तयार होतो. $A$ टप्पा~$S$ वर तयार झाला असल्याने $A$ चे सर्व सदस्य टप्पा~$S$ च्या आधी तयार झालेले असतात (\stagesacc{} नुसार). आता $A$ चे सदस्य असलेले आणि $\phi$ पूर्ण करणारे सर्व संच घ्या; हेसुद्धा टप्पा~$S$ च्या आधीच तयार झालेले असतात. म्हणून \stagesacc{} नुसार टप्पा~$S$ वर त्यांचा $\Setabs{x \in A}{\phi(x)}$ हा संच तयार होतो.

अनौपचारिक गुणधर्माधारित संचरचनेच्या विपरीत, विभक्तीकरणामुळे रसेलची विरोधापत्ती टळते. कारण $\Setabs{x}{x \notin
x}$ हा संच आहे, असे थेट म्हणता येत नाही. पण एखादा संच~$A$ \emph{दिला} असेल, तर $R_A = \Setabs{x \in A}{x \notin x}$ हा संच आहे, असे म्हणता येते. यातून फक्त $R_A \notin R_A$ आणि $R_A \notin A$ हे निष्पन्न होते; त्यांत चिंतेचे काही नाही.

मात्र विभक्तीकरणाचा एक तात्काळ आणि महत्त्वाचा परिणाम आहे:

\begin{thm}\label{sth:z:sep:thm:NoUniversalSet}
\emph{सर्वव्यापी} संच अस्तित्वात नाही; म्हणजे $\Setabs{x}{x = x}$ हा संच अस्तित्वात नाही.
\end{thm}

\begin{proof}
व्यतिरेकासाठी $V$ हा सर्वव्यापी संच आहे असे मानू. मग विभक्तीकरणानुसार $R =
\Setabs{x \in V}{x \notin x} = \Setabs{x}{x \notin x}$ हा संच अस्तित्वात असतो; पण याने रसेलची विरोधापत्ती निर्माण होते.
\end{proof}

सर्वव्यापी संच नसणे---खरे तर संचांच्या पदानुक्रमाला शेवट नसणे---ही संचांच्या संचयी-क्रमिक संकल्पनेमागील एक मूलभूत कल्पना आहे. त्यामुळे सहजज्ञानाने याच निष्कर्षापर्यंत दुसऱ्या मार्गाने कसे पोहोचता येते, हे पाहणे उपयुक्त ठरेल. सर्वव्यापी संच स्वतःचा घटक असलाच पाहिजे. पण टप्पे सुक्रमित असल्याच्या गृहीतकानुसार प्रत्येक संच प्रथम तयार होण्याचा एक टप्पा असतो. त्या टप्प्यावर तयार होणाऱ्या संचाचे सर्व घटक आधीच्या टप्प्यांवर तयार झालेले असतात. म्हणून \emph{कोणताही} संच स्वतःचा सदस्य नसतो: तो स्वतःचा सदस्य असता, तर प्रथम तयार होण्याच्या टप्प्याच्या आधीच तो तयार झाला असता; हा व्याघात आहे. या युक्तिवादाला सर्व सदस्यांचा एखादा शेवटचा टप्पा किंवा त्याचा लगतचा उत्तरवर्ती टप्पा आवश्यक नाही. (संचाच्या ``क्रमस्थानाची'' संकल्पना वापरून हा युक्तिवाद अधिक काटेकोर कसा करायचा, ते \ref{sth:spine:rank:defnsetrank} मध्ये पाहू. मात्र त्यासाठी आधी आणखी काही स्वयंसिद्धके लागतील.)

विभक्तीकरण आणि विस्तारात्मकता यांचे आणखी काही परिणाम पाहू.

\begin{prop}\label{sth:z:sep:prop:emptyexists}
कोणताही संच अस्तित्वात असेल, तर $\emptyset$ अस्तित्वात असतो.
\end{prop}

\begin{proof}
$A$ हा संच असेल, तर विभक्तीकरणानुसार $\emptyset = \Setabs{x \in A}{x \neq x}$ अस्तित्वात असतो.
\end{proof}

\begin{prop}
कोणतेही संच $A$ आणि $B$ यांच्यासाठी $A \setminus B$ अस्तित्वात असतो.
\end{prop}

\begin{proof}
विभक्तीकरणानुसार $A \setminus B = \Setabs{x \in A}{x \notin B}$ अस्तित्वात असतो.
\end{proof}

याशिवाय (जवळजवळ) मनमानी छेदही अस्तित्वात असतात:

\begin{prop}\label{sth:z:sep:prop:intersectionsexist}
$A \neq \emptyset$ असेल, तर $\bigcap A = \Setabs{x}{(\forall y \in A)x \in y}$ अस्तित्वात असतो.
\end{prop}

\begin{proof}
$A \neq \emptyset$ मानू; मग एखादा $c \in A$ असतो. म्हणून $\bigcap A =
\Setabs{x}{(\forall y \in A)x \in y} = \Setabs{x \in c}{(\forall y \in
A)x \in y}$; हा संच विभक्तीकरणानुसार अस्तित्वात असतो.
\end{proof}

पण $A \neq \emptyset$ ही अट लक्षात घ्या. कारण $\bigcap
\emptyset$ हा रिक्ततेमुळे सर्वव्यापी संच ठरेल; ते \ref{sth:z:sep:thm:NoUniversalSet} शी विसंगत आहे.







\section{संयोग}\label{sth:z:union:sec}

\ref{sth:z:sep:prop:intersectionsexist} मधून आपल्याला छेद मिळाले. पण मनमानी संयोगही अस्तित्वात असावेत, तर आणखी एक स्वयंसिद्धक मांडावे लागेल:

\begin{axiom}[संयोग] कोणत्याही संच $A$ साठी $\bigcup A =
\Setabs{x}{(\exists b \in A) x \in b}$ हा संच अस्तित्वात असतो.
\[
	\forall A \exists U \forall x(x \in U \liff (\exists b \in A)x \in b)
\]
\end{axiom}

संचांची संचयी-क्रमिक संकल्पनाही या स्वयंसिद्धकाचे समर्थन करते. $A$ हा संच मानू; \stageshier{} नुसार $A$ कोणत्या तरी टप्प्या $S$ वर तयार होतो. \stagesacc{} नुसार $A$ चा प्रत्येक सदस्य $S$ च्या \emph{आधी} तयार झालेला असतो; त्याचप्रमाणे विचार केल्यास $A$ च्या प्रत्येक सदस्याचा प्रत्येक सदस्यही $S$ च्या आधी तयार झालेला असतो. त्यामुळे \emph{हे सर्व} संच $S$ च्या आधी उपलब्ध असतात आणि $S$ वर त्यांचा एक संच तयार होतो. तोच $\bigcup A$ आहे.







\section{जोड्या}\label{sth:z:pairs:sec}

आता पुढील स्वयंसिद्धक पाहू:

\begin{axiom}[जोड्या]
कोणतेही संच $a, b$ दिल्यास $\{a, b\}$ हा संच अस्तित्वात असतो.
\[
	\forall a \forall b \exists P \forall x (x \in P \liff (x = a \lor x = b))
\]
\end{axiom}

क्रमिक संकल्पनेच्या आधारे या स्वयंसिद्धकाचे समर्थन असे करता येईल. $a$ टप्पा $S$ वर आणि $b$ टप्पा $T$ वर उपलब्ध आहेत असे मानू. $S$ आणि $T$ यांपैकी नंतर येणारा टप्पा $M$ घ्या; दोन्ही टप्पे समान असतील, तर $M=S=T$ घ्या. मग $a$ आणि $b$ दोन्ही टप्पा $M$ वर उपलब्ध असल्याने, $\{a,b\}$ हा $M$ नंतरच्या कोणत्याही टप्प्यावर उपलब्ध वस्तूंपासून बनू शकणारा संग्रह आहे.

पण थांबा!{} $M$ नंतर टप्पे \emph{आहेतच} असे का मानावे? तसे नसतील तर आपले समर्थन अपुरे पडेल. म्हणून जोड्यांचे स्वयंसिद्धक समर्थित करण्यासाठी \ref{sth:z:story:sec} मधील मांडणीत आणखी एक तत्त्व घालावे लागेल:
\begin{enumerate}
\item[] \stagessucc. शेवटचा टप्पा नसतो.
\end{enumerate}
या तत्त्वाचे समर्थन करता येते का? शोनफिल्ड यांच्या मांडणीत शेवटचा टप्पा नाही असे \emph{स्पष्टपणे} म्हटलेले नाही. तरी काटेकोरपणे पाहता हा अतिरिक्त भाग असला, तरी टप्प्याटप्प्याने संच तयार होतात या मूळ कल्पनेशी तो सुसंगत आहे. पुढील चर्चेत आपण हे तत्त्व स्वीकारू. त्याबरोबर जोड्यांचे स्वयंसिद्धकही स्वीकारू.

या नव्या स्वयंसिद्धकाच्या साहाय्याने आणखी अनेक संचांचे अस्तित्व सिद्ध करता येते. उदाहरणार्थ:

\begin{prop}\label{sth:z:pairs:prop:pairsconsequences}
कोणतेही संच $a$ आणि $b$ दिल्यास पुढील संच अस्तित्वात असतात:
\begin{enumerate}
\item\label{sth:z:pairs:singleton} $\{a\}$
\item\label{sth:z:pairs:binunion} $a \cup b$
\item\label{sth:z:pairs:tuples} $\tuple{a, b}$
\end{enumerate}
\end{prop}

\begin{proof}
\ref{sth:z:pairs:singleton}. जोड्यांच्या स्वयंसिद्धकानुसार $\{a, a\}$ अस्तित्वात असतो; विस्तारात्मकतेनुसार तो $\{a\}$ च असतो.

\ref{sth:z:pairs:binunion}. जोड्यांच्या स्वयंसिद्धकानुसार $\{a, b\}$ अस्तित्वात असतो. आता संयोगाच्या स्वयंसिद्धकानुसार $a \cup b = \bigcup
\{a, b\}$ अस्तित्वात असतो.

\ref{sth:z:pairs:tuples}. \ref{sth:z:pairs:singleton} नुसार $\{a\}$ अस्तित्वात असतो. जोड्यांच्या स्वयंसिद्धकानुसार $\{a,
b\}$ अस्तित्वात असतो. मग पुन्हा जोड्यांच्या स्वयंसिद्धकानुसार $\{\{a\}, \{a, b\}\} = \tuple{a, b}$ अस्तित्वात असतो.
\end{proof}

\begin{prob}
कोणतेही संच $a, b, c$ दिल्यास $\{a, b, c\}$ अस्तित्वात असतो, हे दाखवा.
\end{prob}

\begin{prob}
कोणतेही संच $a_1, \ldots, a_n$ दिल्यास $\{a_1, \ldots,
a_n\}$ अस्तित्वात असतो, हे दाखवा.
\end{prob}












\section{घातसंच}\label{sth:z:power:sec}

आता आणखी एक स्वयंसिद्धक पाहू:

\begin{axiom}[घातसंच]
कोणत्याही संच $A$ साठी $\Pow{A} = \Setabs{x}{x \subseteq A}$ हा संच अस्तित्वात असतो.
\[
	\forall A \exists P \forall x(x \in P \liff (\forall z \in x)z \in A)
\]
\end{axiom}

याचे समर्थन अगदी सरळ आहे. $A$ हा संच टप्पा $S$ वर तयार झाला असे मानू. मग \stagesacc{} नुसार $A$ चे सर्व सदस्य $S$ च्या आधी उपलब्ध होते. म्हणून विभक्तीकरणाच्या समर्थनाप्रमाणे विचार केल्यास $A$ चा प्रत्येक उपसंच टप्पा $S$ पर्यंत तयार झालेला असतो. ते सर्व उपलब्ध असल्याने $S$ नंतरच्या कोणत्याही टप्प्यावर त्यांचा एक संच तयार करता येतो. \stagessucc{} नुसार $S$ हा शेवटचा टप्पा नाही, म्हणून असा पुढील टप्पा आहेच. त्यामुळे $\Pow{A}$ अस्तित्वात असतो.

घातसंचाच्या स्वयंसिद्धकाचा एक चांगला परिणाम असा:

\begin{prop}\label{thm:Products}
कोणतेही संच $A, B$ दिल्यास त्यांचा कार्तीय गुणाकार $A \times B$ अस्तित्वात असतो.
\end{prop}

\begin{proof}
घातसंचाचे स्वयंसिद्धक आणि \ref{sth:z:pairs:prop:pairsconsequences} यांनुसार $\Pow{\Pow{A \cup B}}$ अस्तित्वात असतो. म्हणून विभक्तीकरणानुसार पुढील संच अस्तित्वात असतो:
\[
	C = \Setabs{z \in \Pow{\Pow{A \cup B}}}{(\exists x \in A)(\exists y \in B) z = \tuple{x, y}}.
\]
आता कोणतेही $x \in A$ आणि $y \in B$ दिल्यास \ref{sth:z:pairs:prop:pairsconsequences} नुसार $\tuple{x, y}$ अस्तित्वात असतो. शिवाय $x, y
\in A \cup B$ असल्याने $\{x\}, \{x, y\} \in \Pow{A \cup B}$ आणि $\tuple{x,y} \in \Pow{\Pow{A \cup B}}$ मिळते. म्हणून $A \times B = C$.
\end{proof}

या सिद्धतेत घातसंच आणि विभक्तीकरण एकत्र काम करतात. त्यात आश्चर्य नाही. विभक्तीकरण नसते, तर घातसंचाचे तत्त्व फारसे \emph{सामर्थ्यवान} ठरले नसते. कारण कोणते उपसंच अस्तित्वात असतात हे विभक्तीकरण सांगते; त्यामुळे प्रत्येक घातसंच किती ``भरलेला'' असेल, हेही त्यावर अवलंबून असते.

\begin{prob}
कोणतेही संच $A, B$ दिल्यास दाखवा: (i) प्रांत $A$ आणि व्याप्ती $B$ असलेल्या सर्व संबंधांचा संच अस्तित्वात असतो; आणि (ii) $A$ पासून $B$ कडे जाणाऱ्या सर्व फलनांचा संच अस्तित्वात असतो.
\end{prob}

\begin{prob}
$A$ हा संच आणि $\sim$ हा $A$ वरील तुल्यता संबंध मानू. $\sim$ नुसार $A$ वरील तुल्यतावर्गांचा संच, म्हणजे $\equivclass{A}{\sim}$, अस्तित्वात असतो हे सिद्ध करा.
\end{prob}








\section{अनंतता}\label{sth:z:infinity-again:sec}

(एखादा संच अस्तित्वात असेल, तर) अनंत संख्येने संच अस्तित्वात आहेत, हे दाखवण्यासाठी आपल्याकडे आधीच पुरेशी स्वयंसिद्धके आहेत. समजा एखादा संच अस्तित्वात आहे. मग \ref{sth:z:sep:prop:emptyexists} नुसार $\emptyset$ अस्तित्वात असतो. आता कोणताही संच~$x$ दिल्यास \ref{sth:z:pairs:prop:pairsconsequences} नुसार $x \cup \{x\}$ अस्तित्वात असतो. ही प्रक्रिया काही वेळा केल्यास पुढील संच मिळतात:
\begin{enumerate}
	\item[0.] $\emptyset$
	\item[1.] $ \{\emptyset\}$
	\item[2.] $\{\emptyset, \{\emptyset\}\}$
	\item[3.] $\{\emptyset, \{\emptyset\}, \{\emptyset, \{\emptyset\}\}\}$
	\item[4.] $\{\emptyset, \{\emptyset\}, \{\emptyset, \{\emptyset\}\}, \{\emptyset, \{\emptyset\}, \{\emptyset, \{\emptyset\}\}\}\}$
\end{enumerate}
हे प्रत्येक संच इतरांहून वेगळे आहेत, हे तपासता येते.

क्रमांकन $0$ पासून सुरू करण्यामागे काही कारणे आहेत. त्यांपैकी एक असे: ``$n$'' असे नाव दिलेल्या संचात नेमके $n$ सदस्य आहेत आणि सहजज्ञानानुसार तो $n$ व्या टप्प्यावर तयार होतो, हे तपासणे अवघड नाही.

पण यातून आपल्याला \emph{अनंत संख्येने} संच मिळतात; \emph{अनंत संच}, म्हणजे अनंत सदस्य असलेला एक संच, अस्तित्वात आहे याची खात्री मिळत नाही. हा फरक महत्त्वाचा आहे. कारण डेडेकिंड-अनंत संच मिळाल्याशिवाय डेडेकिंड बीजसंरचना बनवता येणार नाही. आणि नैसर्गिक संख्यांचा संच म्हणून वापरता यावी, यासाठी अशी बीजसंरचना आपल्याला हवी आहे. (\ref{sfr:infinite:dedekindsproof:sec} शी तुलना करा.)

आतापर्यंत मांडलेली स्वयंसिद्धके कोणत्याही अनंत संचाचे अस्तित्व \emph{सिद्ध करत नाहीत}. म्हणून नवे स्वयंसिद्धक मांडावे लागेल:

\begin{axiom}[अनंतता]
असा एक संच $I$ अस्तित्वात असतो की $\emptyset \in I$ आणि $x \in I$ असेल तेव्हा $x \cup \{x\} \in I$.
\begin{align*}
	\exists I( & (\exists o \in I)\forall x\ x \notin o \land {}\\
	& (\forall x \in I)(\exists s \in I)\forall z(z \in s \liff (z \in x \lor z = x)))
\end{align*}
\end{axiom}

अनंततेच्या स्वयंसिद्धकातील $I$ साठी $s(x)=x\cup\{x\}$ हे
$I$ पासून $I$ कडे जाणारे फलन आहे. मात्र सर्व $I$ वर त्याचा
एकास-एकपणा थेट गृहीत धरण्याऐवजी प्रथम संवरण-संच घेऊ:
\begin{defn}\label{sth:z:infinity-again:defnomega}
अनंततेच्या स्वयंसिद्धकातून मिळालेला कोणताही संच $I$ घ्या.
$s$ हे $s(x)=x\cup\{x\}$ असे फलन आणि
$\omega=\closureofunder{s}{\emptyset}$ हा $I$ मधील सर्वांत लहान
$s$-बंद, $\emptyset$ धारण करणारा उपसंच असू द्या.
$\omega$ चे सदस्य \emph{नैसर्गिक संख्या} म्हणू.
$n$ ही $\emptyset$ वर $s$ च्या $n$ वेळा केलेल्या उपयोगाचे फलित आहे,
असे म्हणू.
\end{defn}
हे संवरण विभक्तीकरणाने अस्तित्वात येते: $I$ मधील,
$\emptyset$ धारण करणाऱ्या प्रत्येक $s$-बंद उपसंचात असणारे सदस्य घ्या.
असा किमान एक उपसंच $I$ स्वतः आहे. मिळणारा संच
$\emptyset$ धारण करतो, $s$-बंद असतो आणि या सर्व उपसंचांत समाविष्ट असतो.
याच लघुत्वामुळे $\omega$ वर विगमन मिळते: एखादा गुणधर्म
$\emptyset$ साठी सत्य आणि $s$ च्या उपयोगाने जतन होत असेल,
तर त्या गुणधर्माचे $\omega$ मधील सदस्य विभक्तीकरणाने
असा बंद उपसंच बनवतात; म्हणून तो संपूर्ण $\omega$ असतो.
यासाठी आधी $s$ एकास-एक आहे असे मानावे लागत नाही.
या विगमनाने प्रत्येक $n\in\omega$ हा संक्रमक संच,
म्हणजे $(\forall u\in n)u\subseteq n$, असून $n\notin n$ हे सिद्ध होते.
रिक्त संचासाठी दोन्ही अटी स्पष्ट आहेत.
संक्रमक $n$ पासून $s(n)$ संक्रमक होतो: त्याचा सदस्य
$n$ असतो किंवा $n$ चा सदस्य असतो; दोन्ही प्रकरणांत
त्या सदस्याचे सर्व सदस्य $s(n)$ मध्ये असतात.
शिवाय $s(n)\in s(n)$ मानल्यास एकतर $s(n)=n$ किंवा $s(n)\in n$.
पहिल्या प्रकरणात $n\in s(n)$ वरून $n\in n$;
दुसऱ्यात $n$ च्या संक्रमकतेमुळे, पुन्हा $n\in s(n)$ वरून $n\in n$.
दोन्ही प्रकरणांत व्याघात आहे.
आता $s(n)=s(m)$ आणि $n\neq m$ असेल, तर
$n\in m$ आणि $m\in n$; $n$ च्या संक्रमकतेने $n\in n$,
हा व्याघात. त्यामुळे $s$ चे $\omega$ वरील बंधन एकास-एक आहे.
तसेच $\emptyset$ कोणत्याही $s(n)$ चे मूल्य नाही,
कारण $n\in s(n)$. म्हणून $\omega$ डेडेकिंड-अनंत आहे,
आणि $(\omega,s,\emptyset)$ ही डेडेकिंड बीजसंरचना आहे;
\ref{sfr:infinite:dedekind:thm:DedekindInfiniteAlgebra}
शी तुलना करा. $I$ स्वतःही डेडेकिंड-अनंत आहे:
$\omega$ वर $s$ आणि $I\setminus\omega$ वरील प्रत्येक घटकाला तोच घटक देणारे फलन वापरल्यास
$I$ वर एकास-एक फलन मिळते, ज्याच्या व्याप्तीत $\emptyset$ नसतो. या दोन भागांची मूल्ये अनुक्रमे $\omega$ आणि
$I\setminus\omega$ मध्ये असल्याने ती एकमेकांत जुळत नाहीत.
या फलनाचा आलेख $I\times I$ मधून विभक्तीकरणाने मिळतो.

आधीच्या काही परिच्छेदांत ``$n$'' अशी खूण केलेला संच आता $n$ ही संख्याच \emph{मानला} जाईल, हे मागे पाहून तपासता येते.

या निर्णयाचे महत्त्व \ref{sth:z:nat:sec} मध्ये पाहू. सध्या त्याच्या आधारे एक सहज वाटणारा निष्कर्ष सिद्ध करता येतो:

\begin{prop}\label{sth:z:infinity-again:naturalnumbersarentinfinite}
कोणतीही नैसर्गिक संख्या डेडेकिंड-अनंत नाही.
\end{prop}

\begin{proof}
सिद्धता विगमनाने करायची आहे; \ref{sfr:infinite:induction:thm:dedinfiniteinduction} पाहा. स्पष्टच, $0
= \emptyset$ डेडेकिंड-अनंत नाही. विगमनाच्या पायरीसाठी आपण व्यत्यस्त रूप सिद्ध करू: $s(n)$ डेडेकिंड-अनंत असल्याचे (अशक्य असूनही) मानल्यास $n$ डेडेकिंड-अनंत ठरेल.

म्हणून $s(n)$ डेडेकिंड-अनंत आहे असे मानू. म्हणजे असा एखादा एकास-एक फलन $f$ आहे की $\ran{f}\subsetneq \dom{f} = s(n) = n \cup
\{n\}$. आता दोन प्रकरणे पाहावी लागतील.

\emph{प्रकरण १: $n \notin \ran{f}$.} मग $\ran{f} \subseteq n$ आणि $f(n) \in n$. $g = \funrestrictionto{f}{n}$ मानू; मग $\ran{g} =
\ran{f} \setminus \{f(n)\} \subsetneq n = \dom{g}$. म्हणून $n$ डेडेकिंड-अनंत आहे.

\emph{प्रकरण २: $n \in \ran{f}$.} एखादा $m \in \dom{f} \setminus \ran{f}$ निश्चित करू आणि प्रांत $s(n) = n \cup \{n\}$ असलेले फलन $h$ पुढीलप्रमाणे ठरवू:
\[
h(x) =
	\begin{cases}
		f(x) & \text{जर }f(x) \neq n\\
		m & \text{जर }f(x)=n
	\end{cases}
\]
म्हणून $h$ आणि $f$ सर्वत्र जुळतात; फक्त $h(f^{-1}(n)) = m
\neq n = f(f^{-1}(n))$ या ठिकाणी वेगळे आहेत. $f$ हे एकास-एक फलन असल्याने $n \notin
\ran{h}$; तसेच $\ran{h} \subsetneq \dom{h} = s(n)$.
हे $h$ एकास-एक आहे: $f$ चे फक्त एक मूल्य बदलले आहे,
आणि नवे मूल्य $m$ आधीच्या व्याप्तीत नव्हते.
आता प्रकरण १ मधील युक्तिवादाने $n$ डेडेकिंड-अनंत ठरतो.
\end{proof}

पण अनंततेच्या स्वयंसिद्धकाचे \emph{समर्थन} कसे करायचे, हा प्रश्न उरतो. थोडक्यात उत्तर असे की आपण सांगत आलेल्या मांडणीत आणखी एक तत्त्व जोडावे लागेल. ते असे:
\begin{enumerate}
\item[] \stagesinf. अनंत टप्पा असतो. म्हणजे असा टप्पा असतो की (a) तो पहिला नसतो, (b) त्याच्या आधी काही टप्पे असतात आणि (c) त्याचा लगतचा पूर्ववर्ती टप्पा नसतो.
\end{enumerate}
या तत्त्वावरून अनंततेचे स्वयंसिद्धक सरळ मिळते. नैसर्गिक संख्या $n$ ही टप्पा $n$ वर तयार होते, हे आपल्याला ठाऊक आहे. म्हणून $\omega$ हा संच पहिल्या अनंत टप्प्यावर तयार होतो. $\omega$ स्वतःच अनंततेच्या स्वयंसिद्धकासाठी अपेक्षित उदाहरण ठरतो.

मात्र यामुळे \stagesinf{} चे समर्थन कसे करायचे हा प्रश्नच पुढे ढकलला जातो. \stagessucc{} प्रमाणे हे तत्त्वही संचांच्या संचयी-क्रमिक पदानुक्रमाबद्दलच्या मूळ मांडणीत स्पष्टपणे नव्हते. त्याहून महत्त्वाचे म्हणजे संच टप्प्याटप्प्याने तयार होतात या कल्पनेतूनच प्रक्रियेत \emph{अनंत} टप्पा असलाच पाहिजे, असे निष्पन्न होत नाही. टप्प्यांचा क्रम नैसर्गिक संख्यांसारखा आहे, असे मानणे पूर्णपणे सुसंगत वाटते.

मग एक उघड अडचण उभी राहते. \ref{sfr:infinite:dedekindsproof:sec} मध्ये विचारांचा डेडेकिंड-अनंत संच आहे या डेडेकिंड यांच्या ``सिद्धतेचा'' आपण विचार केला. ती तुम्हाला फार समाधानकारक वाटली नसेल. पण संचांच्या क्रमिक संकल्पनेतून (किंवा ``विचारांच्या नियमां''तून) \stagesinf{} आपोआप स्वीकारणे भाग पडत नसेल, तर डेडेकिंड-अनंत संच अस्तित्वात आहे या दाव्याचे अंतर्गत समर्थन अजूनही मिळत नाही.

या विषयावर आणखी बरेच काही सांगता येईल. पण एखादे स्वयंसिद्धक किंवा तत्त्व समर्थित करण्यासाठी \emph{काय लागते}, याचा विचार सुरू करता येईल अशा टप्प्यावर तुम्ही आता पोहोचला असाल, अशी आशा आहे. पुढे आपण \stagesinf{} गृहीत धरू.







\section{$\Z^-$: एक महत्त्वाचा टप्पा}\label{sth:z:milestone:sec}

पुढील विभागात \stagesinf{} कडे पुन्हा वळू. पण अनंततेचे स्वयंसिद्धक मिळाल्यावर आपण एक महत्त्वाचा टप्पा गाठला आहे. आता $\Zminus$ या उपपत्तीसाठी लागणारी सर्व स्वयंसिद्धके आपल्याकडे आहेत. नेमके सांगायचे तर:

\begin{defn}
$\Zminus$ या उपपत्तीची स्वयंसिद्धके अशी आहेत: विस्तारात्मकता, संयोग, जोड्या, घातसंच, अनंतता आणि विभक्तीकरण-रूपबंधाची सर्व उदाहरणे.
\end{defn}

या नावातील \emph{झर्मेलो} म्हणजे झर्मेलो यांची संच उपपत्ती; \emph{वजा} केलेली गोष्ट पुढे पाहू. झर्मेलो यांना हे श्रेय उचित आहे, कारण त्यांनी 1908 मधील कामात ही उपपत्ती मूलत: मांडली होती.\footnote{इतिहास आणि तांत्रिक तपशीलांवरील रंजक टिप्पण्यांसाठी पॉटर (2004, परिशिष्ट A) पाहा.}

या उपपत्तीत गणिताचा फार मोठा भाग मांडता येतो. विशेषत: \ref{sfr:::part} पुन्हा पाहून, तेथे अनौपचारिकरीत्या केलेल्या प्रत्येक गोष्टीची अधिक औपचारिक मांडणी $\Zminus$ मध्ये करता येते, याची तुम्ही \emph{खात्री करून घ्यायला हवी}. (थोडे असे केल्यावर पुढील \ref{sth:z:arbintersections:sec} कडे थेट वळावेसे वाटू शकते.) म्हणून पुढे वेगळी नोंद न करता आपण किमान $\Zminus$ मध्ये काम करत आहोत, असे मानू.







\section{आपल्या नैसर्गिक संख्यांची निवड}\label{sth:z:nat:sec}





\ref{sth:z:infinity-again:defnomega} मध्ये आपण ``नैसर्गिक संख्या'' या शब्दप्रयोगाची स्पष्ट व्याख्या केली. हा निर्णय कसा समजायचा? तो आधिभौतिक दावा नाही; काही संचांना नैसर्गिक संख्या म्हणून \emph{वागवण्याचा} तो निर्णय आहे. असे मानण्याची कारणे आपण \ref{sfr:rel:ref:sec}, \ref{sfr:arith:ref:sec} आणि \ref{sfr:infinite:dedekindsproof:sec} मध्ये पाहिली. पण ती आणखी नेमकी सांगता येतील.

आपले अनंततेचे स्वयंसिद्धक फोन न्यूमन (1925) यांच्या पद्धतीनुसार आहे. मात्र त्याऐवजी पुढील स्वयंसिद्धकही स्वीकारता आले असते:

\begin{defish}
\emph{झर्मेलो यांचे 1908 मधील अनंततेचे स्वयंसिद्धक.} असा एक संच $A$ आहे की $\emptyset \in A$ आणि $(\forall x \in A)\{x\} \in A$.
\end{defish}

आपण फोन न्यूमन यांच्या धर्तीवरचे अनंततेचे स्वयंसिद्धक न घेता झर्मेलो यांचे घेतले असते, तरी डेडेकिंड-अनंत संच आणि म्हणून डेडेकिंड बीजसंरचना मिळाली असती. झर्मेलो यांच्या पद्धतीत आपल्या बीजसंरचनेचा निवडक घटक पुन्हा $\emptyset$ हाच असता ($0$ चे आपले प्रतिरूप); पण एकास-एक फलन $x \mapsto x \cup \{x\}$ ऐवजी $x \mapsto \{x\}$ असे असते. याचा सर्वांत सरळ परिणाम असा की झर्मेलो यांच्या पद्धतीत $2$ म्हणजे $\{\{\emptyset\}\}$, तर आपण फोन न्यूमन यांच्यासह $2$ म्हणजे $\{\emptyset, \{\emptyset\}\}$ मानतो.

अनंततेच्या या दोन स्वयंसिद्धकांपैकी एकच का निवडावे? मुख्य व्यावहारिक कारण असे की फोन न्यूमन यांची पद्धत अनंतपार संख्यांसाठी चांगली विस्तारते. \ref{sth:ordinals:chap} पासून हे पाहू. मात्र केवळ \emph{अंकगणित करण्याच्या} दृष्टीने दोन्ही पद्धती तितक्याच उपयोगी आहेत. म्हणून कोणी नैसर्गिक संख्या \emph{संचच आहेत} असे सांगितले, तर लगेच प्रश्न पडतो: \emph{त्या नेमक्या कोणत्या संच आहेत?}

हा नेमका प्रश्न Benacerraf (1965) यांनी प्रसिद्ध केला. पण आपण यापूर्वी अनेकदा पाहिलेल्या व्यापक बाबीचे ते केवळ सर्वांत प्रसिद्ध उदाहरण आहे. मूलभूत मुद्दा असा: संच उपपत्तीच्या साहाय्याने अनेक ``सहजगत्या'' मानल्या जाणाऱ्या प्रकारच्या वस्तूंचे \emph{अनुकरण} करता येते: वास्तव संख्या, परिमेय संख्या, पूर्णांक आणि नैसर्गिक संख्या; तसेच क्रमित जोड्या, फलने आणि संबंध. पण संच उपपत्ती अशा अनुकरणाची \emph{एकमेव} निवड देत नाही. त्याच कामाला सरळपणे तितकेच उपयोगी पडले असते असे पर्याय \emph{नेहमीच} असतात.







\section{परिशिष्ट: संवरण, गुणधर्माधारित संचरचना आणि छेद}\label{sth:z:arbintersections:sec}

\ref{sth:z:milestone:sec} मध्ये आपण \ref{sfr:::part} मधील अनौपचारिक काम मागे पाहून ते~$\Zminus$ मध्ये करता येते का हे तपासायला सांगितले. तसे केले असल्यास \emph{एक} मुद्दा अडला असेल: डेडेकिंड यांच्या \emph{संवरणांच्या} मांडणीत वापरलेला \emph{छेद}.

\ref{sfr:infinite:dedekind:Closure} आठवा:
\begin{align*}
  \closureofunder{f}{o} & = \bigcap\Setabs{X}{o \in X \text{ आणि $X$ हा
$f$-बंद आहे}}.
\intertext{याचे सर्वसाधारण रूप पुढीलप्रमाणे असते:}
  C & = \bigcap\Setabs{X}{\phi(X)}.
\end{align*}
पण यामुळे सावध व्हायला हवे. अनौपचारिक गुणधर्माधारित संचरचना अपयशी ठरते, म्हणून $\Setabs{X}{\phi(X)}$ अस्तित्वात आहे याची हमी नाही. मग अशा व्याख्यांत काही \emph{अवैध शॉर्टकट} घेतला जातो आहे असे वाटते.

सुदैवाने तसे नाही; किंवा सध्याच्या रूपात त्या व्याख्या \emph{तशा असल्या} तरी थोडे काटेकोर काम करून त्या वैध करता येतात. \ref{sth:z:sep:prop:intersectionsexist} मध्ये $A \neq \emptyset$ असताना $\bigcap A$ का अस्तित्वात असतो हे सांगताना या कामाची पूर्वसूचना दिली होती. आता ते स्पष्टपणे पाहू.

विस्तारात्मकता मानल्यास $C$ ची $\bigcap\Setabs{X}{\phi(X)}$ अशी व्याख्या करताना आपल्याला खरे तर पुढील अट पूर्ण करणारी वस्तू $C$ हवी असते:
\begin{equation}\label{sth:z:arbintersections:bicondelimarbintersection}
	\forall x(x \in C \liff \forall X(\phi(X) \lif x \in X))\tag{*}
\end{equation}
आता असा \emph{किमान एक} संच $S$ आहे की $\phi(S)$, असे मानू. मग \ref{sth:z:arbintersections:bicondelimarbintersection} मिळवण्यासाठी \emph{विभक्तीकरण} वापरून $C$ ची पुढीलप्रमाणे व्याख्या करता येते:
\[
	C = \Setabs{x \in S}{\forall X(\phi(X) \lif x \in X)}.
\]
या व्याख्येवरून अपेक्षेप्रमाणे \ref{sth:z:arbintersections:bicondelimarbintersection} मिळते, हे तपासण्याचा सराव करा.






याच सर्वसाधारण पद्धतीने छेदाच्या व्याख्यांत अनौपचारिक गुणधर्माधारित संचरचनेचा दिसणारा वापर टाळता येईल. ज्या उदाहरणामुळे ही चर्चा सुरू झाली, त्या $\closureofunder{f}{o}$ बाबतीत हे कसे करायचे ते पाहू. \ref{sfr:infinite:dedekind:closureproperties} च्या सिद्धतेच्या सुरुवातीला $o \in \ran{f}\cup\{o\}$ आणि $\ran{f} \cup \{o\}$ हा $f$-बंद संच आहे, असे आपण पाहिले. म्हणून अपेक्षित संचाची व्याख्या अशी करता येते:
\[
	\closureofunder{f}{o} = \Setabs{x \in \ran{f} \cup \{o\}}{(\forall X \ni o)(X \text{ हा $f$-बंद आहे} \lif x \in X)}.
\]



